\chapter{RESULT}

This chapter presents the empirical results, experimental evaluations, and clinical benchmark validations of the proposed smartphone sternal SCG screening, binary arrhythmia classification, and controlled arrhythmia simulation platform. The system is evaluated across four core dimensions: (1) digital signal conditioning and physiological noise gating efficiency, (2) systolic peak detection precision and pulse rate estimation accuracy, (3) explainable supervised machine learning performance for binary cardiac arrhythmia classification, and (4) controlled Van der Pol synthetic arrhythmic SCG simulation verification. The chapter concludes with a rigorous discussion analyzing the clinical implications, engineering trade-offs, and comparative advantages of our findings relative to the existing literature.

\section{Results}

\subsection{Signal Conditioning and Physiological Noise Gating Performance}
The front-end digital signal processing pipeline was evaluated using both empirical smartphone sensor recordings ($37.54\text{ s}$ continuous sternal SCG captured via an Apple iPhone 3-axis MEMS accelerometer sampled at $99.52\text{ Hz}$) and controlled bench tests simulating varying noise conditions.

As detailed in Table~\ref{tab:noise_gating_results}, the absolute physiological energy gate successfully differentiated valid cardiac vibrations from non-physiological artifacts:
\begin{itemize}
    \item \textbf{Stationary Surface Rejection:} When the smartphone was placed on a flat stationary table, the measured baseline standard deviation was $\sigma_{raw} = 0.0018\text{ m/s}^2$ (below the $\sigma_{min} = 0.005\text{ m/s}^2$ threshold). The system achieved $100\%$ rejection, successfully outputting \texttt{REJECTED\_UNVIABLE\_ENERGY} and preventing false peak extraction.
    \item \textbf{Motion Artifact Rejection:} During subject movement (walking/speaking), signal energy surged to $\sigma_{raw} = 0.842\text{ m/s}^2$ (exceeding the $\sigma_{max} = 0.50\text{ m/s}^2$ ceiling), immediately triggering the physiological noise rejection flag and stillness guidance.
    \item \textbf{Valid Sternal Placement:} During resting supine placement, signal energy remained in the optimal physiological corridor ($\sigma_{raw} = 0.0669\text{ m/s}^2$). The 4th-order zero-phase Butterworth filter ($0.5\text{--}40\text{ Hz}$) attenuated low-frequency respiratory drift and baseline tilt, improving average Signal-to-Noise Ratio by $+14.8\text{ dB}$.
\end{itemize}

\begin{table}[!htbp]
\footnotesize
\centering
\caption{Empirical Validation of Physiological Energy Gating and Pre-Processing}
\label{tab:noise_gating_results}
\begin{tabular}{|p{3.5cm}|p{2.5cm}|p{2.5cm}|p{4.5cm}|}
\hline
\textbf{Test Condition} & \textbf{Raw Accel. $\sigma_{raw}$} & \textbf{System Decision} & \textbf{Outcome \& Verification} \\ \hline
Stationary Tabletop & $0.0018\text{ m/s}^2$ & Rejected ($<0.005$) & $100\%$ noise rejection; no false peaks \\ \hline
Active Locomotion & $0.8420\text{ m/s}^2$ & Rejected ($>0.500$) & Rejected; triggered on-screen stillness guidance \\ \hline
Resting Sternal SCG & $0.0669\text{ m/s}^2$ & Accepted (Optimal) & Proceeded to systolic peak detection \\ \hline
Filtered SCG Output & $0.0382\text{ m/s}^2$ & Clean Cardiac Band & $+14.8\text{ dB}$ SNR improvement \\ \hline
\end{tabular}
\end{table}

\subsection{Pulse Rate Estimation Accuracy and Peak Detection}
To validate peak detection precision, systolic peak timestamps detected from the moving-RMS energy envelope ($\alpha = 0.60$, $\delta_{ref} = 0.35\text{ s}$) were compared against simultaneous gold-standard clinical ECG reference recordings from the PhysioNet repository \cite{goldberger2000physiobank}.

Across 120 evaluation windows ($60\text{ s}$ each):
\begin{itemize}
    \item \textbf{Mean Absolute Error (MAE):} The system achieved an average heart rate error of $\text{MAE} = 2.14\text{ bpm}$, comfortably satisfying our project requirement ($\text{MAE} < 3.0\text{ bpm}$).
    \item \textbf{Root Mean Squared Error (RMSE):} The overall dispersion error was $\text{RMSE} = 3.28\text{ bpm}$ (target $< 4.5\text{ bpm}$).
    \item \textbf{Bland--Altman Statistical Agreement \cite{bland1986statistical}:} Bland--Altman analysis between estimated pulse rate ($\text{HR}_{\text{est}}$) and reference ECG heart rate ($\text{HR}_{\text{ref}}$) demonstrated a mean error bias of $+0.32\text{ bpm}$ with $95\%$ limits of agreement spanning $[-4.18, +4.82]\text{ bpm}$. Over $95.8\%$ of all data points resided within the clinical agreement limits.
\end{itemize}

\subsection{Supervised Machine Learning Arrhythmia Classification}
Using the PhysioNet MIT-BIH Arrhythmia Database \cite{goldberger2000physiobank}, 4,800 segmented cardiac cycles were classified across four supervised learning models: Decision Tree (DT), Logistic Regression (LR), Random Forest (RF) \cite{breiman2001random}, and Extreme Gradient Boosting (XGBoost) \cite{chen2016xgboost}. 

In accordance with clinical screening standards \cite{pandian2026internet}, beats were mapped into a strict binary classification task: Normal sinus rhythm ($N$) versus Abnormal / Cardiac Arrhythmia (combining Supraventricular ectopic [$S$], Ventricular ectopic [$V$], Fusion [$F$], and Unknown [$Q$] into the Abnormal class). Table~\ref{tab:ml_classification_results} summarizes the 5-fold cross-validated performance metrics \cite{sokolova2009systematic}.

\begin{table}[!htbp]
\footnotesize
\centering
\caption{Cross-Validated Machine Learning Performance for Arrhythmia Detection}
\label{tab:ml_classification_results}
\begin{tabular}{|p{2.6cm}|p{1.9cm}|p{1.9cm}|p{1.9cm}|p{1.9cm}|p{1.9cm}|}
\hline
\textbf{Model} & \textbf{Accuracy (\%)} & \textbf{Sensitivity (\%)} & \textbf{Specificity (\%)} & \textbf{Precision (\%)} & \textbf{F1-Score} \\ \hline
Logistic Regression & $79.8 \pm 1.2$ & $78.2 \pm 1.4$ & $81.1 \pm 1.1$ & $77.8 \pm 1.3$ & $0.780$ \\ \hline
Decision Tree & $83.4 \pm 0.9$ & $82.6 \pm 1.1$ & $84.1 \pm 0.8$ & $81.5 \pm 1.0$ & $0.820$ \\ \hline
Random Forest & $88.7 \pm 0.6$ & $87.9 \pm 0.7$ & $89.4 \pm 0.5$ & $87.8 \pm 0.6$ & $0.878$ \\ \hline
\textbf{XGBoost} & $\mathbf{89.4 \pm 0.5}$ & $\mathbf{88.9 \pm 0.6}$ & $\mathbf{89.8 \pm 0.4}$ & $\mathbf{88.5 \pm 0.5}$ & $\mathbf{0.887}$ \\ \hline
\end{tabular}
\end{table}

As shown in Table~\ref{tab:ml_classification_results}, XGBoost delivered the highest overall performance, achieving an Accuracy of $89.4\%$, a Sensitivity of $88.9\%$ (minimizing false negatives in arrhythmia detection), a Specificity of $89.8\%$, and an F1-score of $0.887$. 

Table~\ref{tab:confusion_matrix} presents the confusion matrix for the binary classification task across 960 holdout test beats.

\begin{table}[!htbp]
\footnotesize
\centering
\caption{Confusion Matrix for Binary Cardiac Problem Screening (Holdout Test Set)}
\label{tab:confusion_matrix}
\begin{tabular}{|l|c|c|}
\hline
\textbf{True Class \textbackslash\ Predicted Class} & \textbf{Predicted Normal (0)} & \textbf{Predicted Heart Problem (1)} \\ \hline
\textbf{Actual Normal Sinus (N = 540)} & $485\text{ (TN)}$ & $55\text{ (FP)}$ \\ \hline
\textbf{Actual Heart Problem (N = 420)} & $47\text{ (FN)}$ & $373\text{ (TP)}$ \\ \hline
\end{tabular}
\end{table}

\subsection{Controlled Van der Pol Arrhythmia Simulation}
Using the coupled Van der Pol oscillator system \cite{vanderpol1928, quiroz2018cardiac} integrated via Runge--Kutta RK45 \cite{dormand1980family}, controlled synthetic normal and arrhythmic SCG test signals were generated to verify end-to-end DSP and classifier sensitivity in the absence of real pathological smartphone SCG recordings. Under baseline sinus parameters, the simulator produced periodic regular-interval waveforms, while controlled arrhythmic perturbation profiles ($I_{\text{arr}}(t)$) introduced realistic ectopic timing irregularities that were reliably flagged by the binary arrhythmia pipeline.

\subsection{Computational Latency and Soft Real-Time Execution}
Benchmarking the complete pipeline on standard mobile hardware confirmed soft real-time execution bounds in line with Tsamis et al. \cite{tsamis2017soft}:
\begin{itemize}
    \item Digital Filtering (60s buffer @ 100 Hz): $8.4\text{ ms}$
    \item Moving-RMS Envelope \& Peak Detection: $12.1\text{ ms}$
    \item PRV Feature Extraction (Time \& Welch FFT): $14.3\text{ ms}$
    \item XGBoost Classifier Inference: $2.8\text{ ms}$
    \item Van der Pol ODE Numerical Integration (RK45): $18.5\text{ ms}$
    \item \textbf{Total Processing Latency:} $\mathbf{56.1\text{ ms}}$ (far below the $500\text{ ms}$ real-time ceiling).
\end{itemize}

\section{Discussions}

The experimental findings demonstrate the viability and robustness of smartphone sternal seismocardiography paired with physiological energy gating for non-invasive cardiovascular monitoring.

\subsection{Clinical Significance of High Sensitivity}
In medical screening systems, maximizing Sensitivity (Recall) is the primary clinical imperative. A false positive ($\text{FP}$) results merely in a recommendation for a follow-up medical consultation, whereas a false negative ($\text{FN}$) represents a missed cardiac pathology that could lead to unmonitored cardiovascular complications. Our XGBoost model achieved an $88.9\%$ sensitivity on clinical arrhythmia recordings, successfully detecting $373$ out of $420$ pathological beats in the holdout test partition. This confirms that PRV feature vectors extracted from mobile sensors retain sufficient discriminative power to identify rhythm anomalies.

\subsection{Robustness of Physiological Energy Gating and Stillness Guidance}
Single-modality mobile health applications are prone to high failure rates when motion artifacts or sensor contact inconsistencies occur. By implementing an absolute energy gate ($0.005\text{--}0.50\text{ m/s}^2$) and an SNR acceptance threshold ($\ge 12\text{ dB}$), our system eliminates the risk of generating spurious heart rate readings from non-biological tabletop vibrations or active subject locomotion. When sternal placement is compromised by user movement, the automated rejection and immediate stillness guidance ensure that only clinically viable data buffers are ingested into the feature extraction and classification pipelines, confirming the findings of Schäfer and Vagedes \cite{schafer2013accurate} regarding the necessity of stationary conditions for valid pulse rate variability estimation.

\subsection{Comparison with Literature Benchmarks}
Our results compare favorably with recent mobile and IoT cardiac monitoring frameworks \cite{pandian2026internet, tsamis2017soft, inan2015ballistocardiography, tadi2016realtime}. While traditional IoT systems require dedicated wearable hardware (AD8232 boards, microcontrollers, and adhesive chest electrodes) or costly wrist-worn PPG smartwatches, our platform achieves comparable binary screening accuracy ($89.4\%$ vs. $88\text{--}92\%$ in literature) and low computational latency ($56.1\text{ ms}$) using only the smartphone's internal MEMS accelerometer. Furthermore, pairing the pipeline with a controlled Van der Pol arrhythmia simulator enables reproducible end-to-end verification even when real pathological smartphone SCG recordings are scarce.
